\chapter{Cross-structure orbital mapping (menu 17)}
\label{ch:menu17}
\menulabel{17}

\section{Purpose}

Keep one active space \emph{consistent} along a path --- a reaction coordinate,
a bond-dissociation curve, a scan.  Active spaces selected structure by
structure drift apart (the active-space-inconsistency error), the total
correlation energy is computed inconsistently, and the relative energies become
erratic.  The protocol maps the localized orbitals of the structures onto each
other through two shape-and-localization criteria --- the orbital kinetic energy
and the shell-wise populations --- and transfers a selection through the map:

\[
|t_{iL} - t_{jK}| / E_\mathrm{h} < \tau_{\mathrm{kin}},
\qquad
\sum_a \bigl| q^a_{iL} - q^a_{jK} \bigr| < \tau_{\mathrm{loc}},
\]

with one threshold $\tau = \tau_{\mathrm{kin}} = \tau_{\mathrm{loc}}$.  The
orbitals that change along the path (the bonds being broken or formed) cannot
be matched in every structure and are collected into a non-matchable block; the
whole block joins the active space together whenever one of its orbitals is
selected anywhere, and the same holds for every class the map built.

\section{What you need}

One \code{orca\_2json} export per structure, \emph{localized} (localized
orbitals are what transfers between structures), each carrying the
\code{T-Matrix} (kinetic criterion) and the \code{S-Matrix}/labels
(shell-population criterion) --- export with a
\code{<basename>.json.conf} of
\code{\{"MOCoefficients": true, "1elIntegrals": ["H", "S", "T", "V"]\}}, and
localize with \code{orca\_loc} first (measured: ORCA's IAO-based methods refuse
the virtual block, so the frozen example localizes the occupied block with
IAO-BOYS and the virtual block with Foster-Boys, both with the random seed
off).  All structures must share the basis set, the atom order and the
orientation.  The manifest JSON schema is in Chapter~\ref{ch:formats}.

\section{Walkthrough}

\lstinputlisting[style=fbkcode, firstline=54, lastline=84,
  title={The menu-17 example session (the frozen N$_2$ scan, three localized
  exports, a selection on one bond orbital)}]
  {examples/expected/m17_mapping.txt}

With an \code{active} block in the manifest, the same report carries the
overlap-preservation scalar of the second source (its Supporting Information
Eq.~(2); the main text's Eqs.~(1)--(3) give the general form) for every
adjacent pair, together with the alignment diagnostic $O_{\min}$ of that
source's main text:

\lstinputlisting[style=fbkcode, firstline=236, lastline=240,
  title={The active-space overlap block of the same run (\code{active} =
  the bond triad; $|\det S_{\mathrm{act}}|$ and $O_{\min}$)}]
  {examples/expected/m17_mapping.txt}

The selection given for one structure ($\{4\}$ of the equilibrium geometry)
implies the consistent space of every structure:

\lstinputlisting[style=fbkcode, firstline=157, lastline=160,
  title={The consistent active space the selection implies (same run)}]
  {examples/expected/m17_mapping.txt}

\section{What you get}

Per structure the descriptor table (occupation, kinetic energy in $E_\mathrm{h}$,
largest shell-wise populations), the tau curve, the classes the map built, the
non-matchable sets, and --- with \code{selections} --- each structure's
consistent active space; plus a report file
(\file{<manifest>.mapping.fbk.md}) with the citations.

\section{How to read the numbers}

\begin{readbox}
\begin{itemize}
  \item \textbf{A class} is a set of orbitals that map onto each other across
    \emph{all} structures, in both directions and as a whole set: degenerate
    orbitals ride together (in the example, the three equivalent N--N bond
    orbitals form a class of nine orbital--structure pairs).  The
    non-matchable block is the changing set: here the two one-centre hybrids,
    whose $s/p$ ratio follows the bond length.
  \item \textbf{The tau curve is the feature, not one number.}  Its plateaus
    are where the non-matchable count is stable (the source reads its own
    example at $\tau \le 0.5$; this one's occupied side steps 21 $\to$ 15
    $\to$ 6 and plateaus from $\tau = 0.375$).  Without an explicit
    \code{tau}, Eq.~(8) of the source is applied per orbital set: the
    threshold minimising the non-matchable count, the largest one among ties
    and capped at 0.5.
  \item \textbf{The populations here are shell-wise L\"owdin populations}
    computed from the export; the source uses IAO ones (its own text licenses
    any orbital-wise population analysis, having chosen IAO for basis-set
    insensitivity).  The population \emph{scale} therefore differs from the
    source's, and $\tau$ is a data-calibrated parameter here --- read the
    curve, never a recalled absolute value.
  \item \textbf{The virtual manifold can be genuinely unmatchable}: diffuse
    orbitals reorganize with geometry (in the example the virtual side never
    drops below 35 non-matchable orbitals).  The report shows that instead of
    hiding it; the occupied side is where the protocol's canonical behaviour
    (cores map, changing bonds do not) is read.
  \item \textbf{Two numbers, one block.}  $|\det S_{\mathrm{act}}|$ multiplies
    the block's singular values, so one poor direction drags it down, while
    $O_{\min}$ \emph{is} the smallest singular value --- the second source's
    alignment diagnostic, whose working criterion ($O_{\min} \ge 0.85$) is
    calibrated on its own condensed-phase dataset (quoted as such in the
    report).  In the example the two 0.5-\AA{} steps read $|\det|$ 0.7172 and
    0.7251 with $O_{\min}$ 0.8567 and 0.8323: only the second crosses the
    source's line, and the block's boundary notes state which setting each
    number belongs to.
\end{itemize}
\end{readbox}

\section{Worked example}

The frozen chain is an N$_2$ scan in RHF/def2-SVP at 1.094, 1.600 and
2.600~\AA{} (\file{examples/run-examples.sh} copies the three localized
exports).  The map reads back the chemistry: the two $1s$ cores map uniquely
across all three structures ($t$ = 22.6155 / 22.5626 / 22.6405~$E_\mathrm{h}$,
agreeing to well within $\tau$), the three equivalent bond orbitals form one
degenerate class, and the two one-centre hybrids --- the valence orbitals whose
character changes most along the stretch --- are the non-matchable block.
Selecting one bond orbital of the equilibrium structure therefore pulls the
whole degenerate triad into every structure's consistent space, and selecting
a hybrid pulls the hybrid block; that is the source's union rule in one line.

\section{Refusals and related sections}

\begin{refusalbox}
\begin{itemize}
  \item An export without the \code{T-Matrix} (or without the
    \code{S-Matrix}/labels) is refused with the re-export recipe; an export
    with fractional occupations (an active space) is refused --- the mapping
    is defined for an occupied/virtual split, so map the Hartree--Fock
    orbitals of the path first and select from the map.
  \item Structures with a different shell set, a different orbital count, or
    fewer than two structures are refused: the mapping compares structures in
    their own frames (same basis, same atom order, same orientation).
  \item The selection itself is not made here: menu~\menu{12}'s entropy
    protocol (or any other) provides it, and this menu reports the consistent
    space it implies.
\end{itemize}
\end{refusalbox}

Related: \menu{13} (the $\sigma_F$/space-change comparison of two spaces of
\emph{one} structure --- this menu is its cross-structure counterpart),
\menu{12} (the entropy selection the \code{selections} block carries), and
Chapter~\ref{ch:criteria} for the criterion table.
